\section{Response geometry}
\label{sec:geometry}

\subsection{The stellar map}

Let \(\eos\) denote a cold barotropic EOS and let \(h_c\) denote the central relativistic enthalpy labeling a stellar configuration. Solving the stellar-structure and perturbation equations defines
\begin{equation}
    \Phi:(\eos,h_c)\longrightarrow \yvec .
    \label{eq:stellar-map}
\end{equation}
The EOS is an infinite-dimensional object. For numerical work we introduce coordinates \(a^i\) on a finite approximation to EOS function space, with the functional limit recovered under basis refinement.

The differential of Eq.~\eqref{eq:stellar-map} separates into an EOS part and a stellar-sequence part,
\begin{equation}
    d y^A
    =
    J^A{}_i\, da^i
    +
    t^A\, dh_c ,
    \label{eq:differential-map}
\end{equation}
where
\begin{equation}
    J^A{}_i
    \equiv
    \left.\frac{\partial y^A}{\partial a^i}\right|_{h_c},
    \qquad
    t^A
    \equiv
    \left.\frac{\partial y^A}{\partial h_c}\right|_{\eos}.
    \label{eq:J-t-defs}
\end{equation}
The vector \(t^A\) is tangent to a fixed-EOS stellar sequence. Quasi-universality concerns the extent to which changing the EOS moves this sequence transversely rather than merely reparameterizing motion along it.

\subsection{Metric-aware EOS response}

A singular-value decomposition of \(J^A{}_i\) is not by itself a coordinate-independent characterization of EOS sensitivity. Under a change of EOS coordinates \(a^i\rightarrow\tilde a^\alpha(a)\), the numerical scale of the columns of \(J\) changes. In particular, a rescaling \(a^i\rightarrow\lambda_i a^i\) changes the singular spectrum without changing the underlying set of EOSs.

We therefore equip EOS perturbations with a positive covariance or inverse metric \(\aeos^{ij}\). In a probabilistic interpretation this can be the local covariance of a nonparametric EOS process; geometrically it defines which perturbations are considered comparable. The EOS-induced second moment in observable space is
\begin{equation}
    \Geos^{AB}
    =
    J^A{}_i\,\aeos^{ij}\,J^B{}_j .
    \label{eq:induced-response}
\end{equation}
Under a consistent coordinate transformation of both \(J\) and \(\aeos\), Eq.~\eqref{eq:induced-response} is invariant.

In the functional limit, choose an EOS field \(s(h)\), for example a latent field controlling the sound speed, and define
\begin{equation}
    K^A(h)
    =
    \frac{\delta y^A}{\delta s(h)} .
    \label{eq:functional-kernel}
\end{equation}
Then
\begin{equation}
    \Geos^{AB}
    =
    \int dh\,dh'\,
    K^A(h)\,
    C_{\rm EOS}(h,h')\,
    K^B(h') .
    \label{eq:functional-G}
\end{equation}
Equation~\eqref{eq:functional-G} is the central first-order object.

\subsection{Quotienting the stellar-sequence direction}

Observable space also requires a metric to define orthogonality and normalize covectors. Let \(g_{AB}\) denote a declared positive metric on observable space. A Euclidean metric in logarithmic observables is a useful baseline, but it is not assumed unique; coordinate dependence is a required robustness test.

The normalized sequence tangent is
\begin{equation}
    \hat t^A
    =
    \frac{t^A}{(g_{BC}t^Bt^C)^{1/2}} .
\end{equation}
The projector onto directions \(g\)-orthogonal to the sequence is
\begin{equation}
    P^A{}_B
    =
    \delta^A{}_B
    -
    \hat t^A \hat t_B ,
    \qquad
    \hat t_B=g_{BC}\hat t^C .
    \label{eq:projector}
\end{equation}
The transverse EOS response is
\begin{equation}
    \Gperp^{AB}
    =
    P^A{}_C\,
    \Geos^{CD}\,
    P^B{}_D .
    \label{eq:Gperp}
\end{equation}

For a two-dimensional observable plane such as \((\ln\Lamtwo,\ln\Ibar)\), the quotient by the one-dimensional sequence tangent leaves a single normal direction. Its response gives a local scalar measure of non-universality. In higher-dimensional observable spaces, the physical hierarchy is obtained from the constrained cotangent-space problem below, equivalently a generalized eigenproblem normalized with \(g^{AB}\); raw Euclidean eigenvalues of \(P G P^T\) are not themselves coordinate invariant.

\subsection{Local quasi-universal covectors}

A covector \(n_A\) defines a local candidate relation when its EOS-induced variance
\begin{equation}
    \mathcal V[n]
    =
    n_A\,\Geos^{AB}\,n_B
    \label{eq:variance-functional}
\end{equation}
is small subject to
\begin{equation}
    n_A t^A=0,
    \qquad
    n_A g^{AB}n_B=1 .
    \label{eq:normalization}
\end{equation}
Equivalently, \(n_A\) is a low-response mode in the normal cotangent space to the stellar sequence. This definition makes prior dependence explicit through \(C_{\rm EOS}\), rather than hiding it in the composition of a finite EOS catalog.

\subsection{From local insensitivity to a global relation}
\label{sec:integrability}

A local normal field is not automatically the gradient of a scalar relation. A global quasi-universal relation requires a function \(F(\yvec)\) such that
\begin{equation}
    F(\yvec)\simeq {\rm const},
    \qquad
    n_A(\yvec)\propto\partial_A F .
    \label{eq:global-relation}
\end{equation}
Thus the existence of a suppressed local response mode and the existence of a global fitted curve or hypersurface are distinct statements.

For a codimension-one distribution represented by the one-form
\begin{equation}
    \omega=n_A\,dy^A,
\end{equation}
Frobenius integrability requires
\begin{equation}
    \omega\wedge d\omega=0
    \label{eq:frobenius}
\end{equation}
where the condition is nontrivial. In the two-observable I--Love plane local line fields are automatically integrable, but path dependence becomes nontrivial when the construction is embedded in a larger observable/configuration space or when several low-response directions are retained. We therefore use both local Frobenius diagnostics in higher dimension and direct path-dependence tests of transported covectors.

Operationally, the global relation is reconstructed after determining \(n_A\) from the stellar response. The standard empirical I--Love fit is used only for comparison.

\subsection{Second order and residual scatter}

Suppose \(F(\Phi)\) is stationary at first order under EOS perturbations. For finite EOS displacement \(\delta a^i\),
\begin{align}
    \delta F
    &=
    \partial_A F\,J^A{}_i\,\delta a^i
    \nonumber\\
    &\quad+
    \frac{1}{2}
    \left[
        \partial_A F\,H^A{}_{ij}
        +
        \partial_A\partial_BF\,
        J^A{}_iJ^B{}_j
    \right]
    \delta a^i\delta a^j
    +\mathcal O(\delta a^3),
    \label{eq:second-order}
\end{align}
with
\begin{equation}
    H^A{}_{ij}
    =
    \frac{\partial^2y^A}{\partial a^i\partial a^j}.
\end{equation}
When the first line of Eq.~\eqref{eq:second-order} is suppressed, the bracketed term is the leading local source of finite variation in the scalar relation.  It contains two conceptually distinct pieces: the curvature \(H^A{}_{ij}\) of the EOS-to-observable map and the curvature \(\partial_A\partial_BF\) of the chosen relation in observable space.  In Sec.~\ref{sec:curvature} we test the first piece directly by freezing the local soft normal and measuring displacement away from its tangent plane; we do not claim there to have separately validated the complete second-order expansion of the globally reconstructed \(F\).

\subsection{Relation to previous stationarity results}

The perturbative results of Refs.~\cite{chan2016stationarity,yip2017tidal} imply a vanishing or strongly suppressed first variation for specific moment--Love combinations near the incompressible limit. In the present notation, such stationarity corresponds schematically to
\begin{equation}
    n_AJ^A{}_i\simeq0
\end{equation}
for the relevant perturbations. The difference here is the order of construction. We do not begin with a chosen \(n_A\). We calculate the EOS response of several observables, quotient the sequence direction, and infer the low-response normal field from the resulting operator. The known I--Love relation is therefore a test of the construction.
